\documentclass[a4paper]{article}

% Shared Preamble for MathLatexDocs

% --- Core Packages ---
\usepackage[english]{babel}
\usepackage[utf8]{inputenc}
\usepackage{amsmath}
\usepackage{amsthm}
\usepackage{amssymb}
\usepackage{mathtools}
\usepackage{graphicx}
\usepackage{microtype} % Better justification and fewer overfull hboxes
\usepackage{bm}
\usepackage{enumitem}
\usepackage{booktabs}
\usepackage{mathrsfs}

% --- Layout & Design ---
\usepackage{docmute}
\usepackage{subfiles} % Support for master/subfile structure
\usepackage[colorinlistoftodos]{todonotes}
\usepackage[colorlinks=true, linkcolor=blue, citecolor=magenta]{hyperref}
\usepackage{cleveref} % Better cross-referencing
\usepackage{tcolorbox} % For styled boxes
\usepackage{tikz} % For drawing diagrams
\usetikzlibrary{arrows.meta,positioning}

% --- Theorem Environments (Classic Style) ---
\makeatletter
\@ifundefined{classicthm}{%
  \newtheorem{classicthm}{Theorem}[section]
  \newtheorem{classicdef}[classicthm]{Definition}
  \newtheorem{classiclem}[classicthm]{Lemma}
  \newtheorem{theorem}{Theorem}[section]
  \newtheorem{corollary}{Corollary}[theorem]
  \newtheorem{lemma}[theorem]{Lemma}
  \newtheorem{definition}{Definition}[section]
  \newtheorem{proposition}[theorem]{Proposition}
  \newtheorem{remark}[theorem]{Remark}
  \newtheorem{example}[theorem]{Example}
}{}

% Wrappers to maintain compatibility with the {Title}{Label} syntax used in files
\@ifundefined{mytheorem}{%
  \newenvironment{mytheorem}[2]{%
    \begin{classicthm}[#1]\label{thm:#2}
  }{%
    \end{classicthm}
  }
  
  \newenvironment{mydefinition}[2]{%
    \begin{classicdef}[#1]\label{def:#2}
  }{%
    \end{classicdef}
  }
  
  \newenvironment{mylemma}[2]{%
    \begin{classiclem}[#1]\label{lem:#2}
  }{%
    \end{classiclem}
  }
}{}
\makeatother

% --- Custom Commands ---
\providecommand{\Z}{\mathbb{Z}}
\providecommand{\R}{\mathbb{R}}
\providecommand{\C}{\mathbb{C}}
\providecommand{\N}{\mathbb{N}}
\providecommand{\G}{\mathcal{G}}
\providecommand{\Lop}{\mathcal{L}}
\providecommand{\PP}{\mathcal{P}}
\providecommand{\dd}{\,\mathrm{d}}
\providecommand{\bitand}{\mathbin{\&}}
\providecommand{\CFDots}{\genfrac{}{}{0pt}{0}{}{\ddots}}


\title{Exact solution to Polynomial equations}
\author{Emil Kerimov}
\date{}


\begin{document}
\maketitle

\section{History}\label{eps:sec:history}

Scipione del Ferro is thought to have been the first to discover the general solution to the cubic equation, but did not publish it, and kept it a secret until his deathbed, telling his student Antonio Fior.
Fior then challenged Niccolo Fontana, also known as Tartaglia, in 1535 to solve a list of 30 depressed cubics.
On the night of February 13, 1535, Tartaglia discovered the solution.

\bigskip

Tartaglia revealed the secret to
Gerolamo Cardano on March 25, 1539, with the oath
of Cardano to never publish the findings.
Cardano, along with his pupil Ludovico Ferrari, discovered the solution of the general cubic equation.

While inspecting del Ferro's papers, they
came across a solution to the depressed cubic in del Ferro's own handwriting.
Cardano and Ferrari were no longer prohibited from publishing their results.

\bigskip

In 1545, Cardano published his book Ars
Magna, the \textquotedblleft Great Art.\textquotedblright
In it, he published the solution to the
depressed cubic, with a preface crediting
del Ferro with the original solution.
He also published his solution of the
general cubic and also Ferrari's solution of
the quartic.

\bigskip

Lagrange published a paper in 1771, Reflections on the Algebraic Theory of Equations, in which he described why the techniques used to solve the quadratic, cubic and quartic equations could not be applied to solve the quintic equation.

\bigskip

Paolo Ruffini attempted to prove that the quintic equation was insolvable.
He published 6 proofs on the subject from 1799 to 1813.
However, other mathematicians found his work too long and difficult to understand, except for Cauchy.
These proofs did end up having a few gaps in them, but were on the right track.

\bigskip

Niels Henrik Abel published his proof that the quintic equation was insolvable in 1824.
Abel is credited for finally putting this question to rest.

\bigskip

Evariste Galois, who died at the age of 20 in 1832 in a duel, studied attempted to generalize how to determine when the roots could be solved in terms of radicals.
He became the first mathematician to use the word 'group' and defined the normal subgroup.

\section{Quadratic}\label{eps:sec:quadratic}

\begin{definition}\label{eps:quad def}
  A quadratic polynomial
  \[
    ax^2 + bx + c = 0
  \]
\end{definition}

Given a quadratic polynomial, the two solutions of $x$ that satisfy this equation are given by the quadratic equation
\begin{equation}
  \boxed{
    x = \frac{-b \pm \sqrt{b^2 - 4ac}}{2a}
  }\label{eps:eq:equation}
\end{equation}

This one equation has three solutions due to the cube root of complex numbers.

\subsection{Algebraic Proof}\label{eps:subsec:algebraic-proof}
The most straight-forward proof is given by a technique referred to as 'Completing the Square'.

\begin{gather*}
  ax^2 + bx + c = 0\\
  \shortintertext{Dividing both sides by $a$}
  x^2 + \frac{b}{a}x + \frac{c}{a} = 0\\
  \shortintertext{Adding and subtracting $\frac{b^2}{(2a)^2}$}
  x^2 + \frac{b}{a}x + \frac{b^2}{(2a)^2} +  \frac{c}{a} - \frac{b^2}{(2a)^2} = 0\\
  \shortintertext{Recognizing the 'square'}
  (x+ \frac{b}{2a})^2 +  \frac{c}{a} - \frac{b^2}{(2a)^2} = 0\\
  \shortintertext{Rearranging and solving}
  \begin{alignedat}{2}
    (x + \frac{b}{2a})^2 & = \frac{b^2}{4a^2}-  \frac{c}{a}  = \frac{b^2-4ac}{4a^2} \\
    x + \frac{b}{2a}     & =\pm \sqrt{ \frac{b^2 - 4ac}{4a^2}}                      \\
    x                    & = - \frac{b}{2a} \pm \sqrt{ \frac{b^2 - 4ac}{4a^2}}      \\
    x                    & = \frac{-b \pm \sqrt{b^2 - 4ac}}{2a}                     \\
  \end{alignedat}
\end{gather*}

\subsection{Discriminant}\label{eps:subsec:discriminant}

The discriminant of a quadratic equation is defined to be
\begin{equation}
  \Delta = b^2 - 4ac\label{eps:eq:equation2}
\end{equation}

\section{Cubic}\label{eps:sec:cubic}

\begin{definition}\label{eps:cubic def}
  A cubic polynomial
  \[
    ax^3 + bx^2 + cx + d = 0
  \]
\end{definition}

Given a cubic polynomial the three solutions of $x$ that satisfy this equation are given by

\begin{equation} \label{eps:eq:cubic-solution}
  \boxed{
    \begin{split}
      x = - \frac{b}{3a} +
      \frac{1}{\sqrt[3]{2} \cdot 3 a} \Bigg(
      \sqrt[3]{9abc - 27 d a^2 - 2 b^3  + 3a \sqrt[2]{
          3 \Big(
          27 a^2 d^2 + 4 b^3 d - 18 abcd + 4 a c^3 - b^2 c^2
          \Big)
        }}
      \\
      +
      \sqrt[3]{9abc - 27 d a^2 - 2 b^3 -
        3a \sqrt[2]{
          3
          \Big(
          27 a^2 d^2 + 4 b^3 d - 18 abcd + 4 a c^3 - b^2 c^2
          \Big)
        }}
      \Bigg)
    \end{split}
  }
\end{equation}

This was first shown by Gerolamo Cardano in 1545

\subsection{Depressed Cubic}\label{eps:subsec:depressed-cubic}
\begin{definition}\label{eps:depressed cubic def}
  A depressed cubic is a polynomial in the form of
  \[
    x^3 + cx + d = 0
  \]
\end{definition}


\begin{theorem}
  We can transform any cubic to a depressed cubic form.

  Proof
  \begin{gather*}
    \intertext{Starting from definition~\ref{eps:cubic def}}
    ax^3 + bx^2 + cx + d = 0
    \\
    \intertext{Dividing by a, we define}
    a_2 = \frac{b}{a}, a_1 = \frac{c}{a}, a_0 = \frac{d}{a}
    \\
    \intertext{Rewriting}
    x^3 + a_2 x^2 + a_1 x + a_0 = 0
    \\
    \intertext{Completing the cube}
    x^3 + 3 \Big( \frac{a_2}{3} \Big) x^2 + 3 \Big( \frac{a_2}{3} \Big)^2 x  + \Big( \frac{a_2}{3} \Big)^3 - 3 \Big( \frac{a_2}{3} \Big)^2 x  - \Big( \frac{a_2}{3} \Big)^3 + a_1 x + a_0 = 0
    \\
    \intertext{Using $(x+t)^3 = x^3 + 3x^2 t + 3x t^2 + t^3$}
    \Big( x + \frac{a_2}{3} \Big)^3 + x \Big( a_1 - 3 \Big( \frac{a_2}{3} \Big)^2 \Big)  + \Big( a_0 - \Big( \frac{a_2}{3} \Big)^3 \Big) = 0
    \\
    \intertext{Let $z=x + \frac{a_2}{3}$}
    z^3 + \Big(z - \frac{a_2}{3} \Big) \Big( a_1 - 3 \Big( \frac{a_2}{3} \Big)^2 \Big)  + \Big( a_0 - \Big( \frac{a_2}{3} \Big)^3 \Big) = 0
    \\
    \intertext{Simplifying and collecting like terms}
    z^3 +
    z \Big( a_1 - 3 \Big( \frac{a_2}{3} \Big)^2 \Big)
    +\Bigg(
    - \frac{a_2}{3} \Big( a_1 - 3 \Big( \frac{a_2}{3} \Big)^2 \Big)
    + \Big( a_0 - \Big( \frac{a_2}{3} \Big)^3 \Big)
    \Bigg) = 0
    \\
    \intertext{Therefore after shifting with $z=x + \frac{b}{3a}$}
    c' = \frac{c}{a} - 3 \Big( \frac{b}{3a} \Big)^2
    \\
    d' = - \frac{b}{3a} \Big( \frac{c}{a} - 3 \Big( \frac{b}{3a} \Big)^2 \Big)
    + \Big( \frac{d}{a} - \Big( \frac{b}{3a} \Big)^3 \Big)
    \intertext{Simplifying}
    c' =  \frac{c}{a} - 3 \Big( \frac{b}{3a} \Big)^2
    \\
    d' =\frac{d}{a} - \Big( \frac{b}{3a} \Big)^3 - \frac{b}{3a} c'
  \end{gather*}
\end{theorem}


\begin{theorem}
  Starting from the depressed cubic form, we can solve for the roots.

  Proof
  \begin{gather*}
    \intertext{Starting from definition~\ref{eps:depressed cubic def}}
    z^3 + cz + d = 0
    \\
    \intertext{Define $p=-\frac{c}{3}$ and $q=-\frac{d}{2} $}
    z^3 = 3pz + 2q
    \\
    \intertext{Consider $z=u+v$}
    z^3 = (u+v)^3 = u^3 + 3u^2 v + 3 u v^2 + v^3 = u^3 + v^3 + 3uv(u+v) = 3uv(z) + u^3 + v^3
    \\
    \intertext{Comparing these forms}
    3uv = 3p \\
    u^3 + v^3 = 2q \\
    \intertext{This implies}
    v = \frac{p}{u}\\
    u^3 + \Big(\frac{p}{u}\Big)^3 = 2q
    \intertext{Therefore we can create a polynomial equation to solve for this relation}
    u^6 + -2q u^3 + p^3 = 0
    \\
    \intertext{We can use the quadratic solution to solve for this relation}
    u^3 = \frac{2q \pm \sqrt{4q^2 - 4p^3}}{2} = q \pm \sqrt{q^2 - p^3}
    \\
    \intertext{Define the solutions for u}
    u_{+} = \sqrt[3]{q + \sqrt{q^2 - p^3}} \\
    u_{-} = \sqrt[3]{q - \sqrt{q^2 - p^3}} \\
    \intertext{Since $u_{+}^3 u_{-}^3  = p^3$, from quadratic equation}
    u_{+}u_{-} = p
    \intertext{Using definition $v = \frac{p}{u}$}
    v_{+} = \frac{p}{u_{+}} =\frac{p u_{-}}{u_{+}u_{-}} = \frac{p u_{-}}{p} =  u_{-}
    \intertext{Similarly}
    v_{-} =  u_{+}
    \intertext{Therefore the solution to $z^3 = 3pz + 2q$}
    z = u + v = u_{+} + u_{-} =  \sqrt[3]{q + \sqrt{q^2 - p^3}} + \sqrt[3]{q - \sqrt{q^2 - p^3}}
    \intertext{Re-writing in the original form}
    z = \sqrt[3]{-\frac{d}{2} + \sqrt{\frac{d^2}{4} +\frac{c^3}{27}}} + \sqrt[3]{-\frac{d}{2} - \sqrt{\frac{d^2}{4} + \frac{c^3}{27}}}
  \end{gather*}
\end{theorem}



\subsection{General Cubic}\label{eps:subsec:general-cubic}
Combining all of these results to back out the solution for a general cubic equation in terms of the original variables, $ax^3 + bx^2 + cx + d = 0$,
we first summarize the steps taken.

\begin{itemize}
  \item shift solution by $z=x + \frac{b}{3a}$
  \item write solution in depressed cubic form
  \item solve the depressed cubic equation.
\end{itemize}

Undoing these substitutions backwards
\begin{gather*}
  z = \sqrt[3]{-\frac{d'}{2} + \sqrt[2]{\frac{d'^2}{4} +\frac{c'^3}{27}}} + \sqrt[3]{-\frac{d'}{2} - \sqrt[2]{\frac{d'^2}{4} + \frac{c'^3}{27}}}
  \\
  \intertext{Using $ c' =  \frac{c}{a} - 3 \Big( \frac{b}{3a} \Big)^2 $
    and
    $d' =\frac{d}{a} - \Big( \frac{b}{3a} \Big)^3 - \frac{b}{3a} c' $}
  z = \sqrt[3]{-\frac{\Big( \frac{d}{a} - \Big( \frac{b}{3a} \Big)^3 - \frac{b}{3a} \Big( \frac{c}{a} - 3 \Big( \frac{b}{3a} \Big)^2 \Big)}{2} + \sqrt[2]{\frac{\Big(\frac{d}{a} - \Big( \frac{b}{3a} \Big)^3 - \frac{b}{3a} \Big( \frac{c}{a} - 3 \Big( \frac{b}{3a} \Big)^2 \Big)^2}{4} +\frac{\Big( \frac{c}{a} - 3 \Big( \frac{b}{3a} \Big)^2 \Big)^3}{27}}}
  \\
  +
  \sqrt[3]{-\frac{\Big( \frac{d}{a} - \Big( \frac{b}{3a} \Big)^3 - \frac{b}{3a} \Big( \frac{c}{a} - 3 \Big( \frac{b}{3a} \Big)^2 \Big)}{2} - \sqrt[2]{\frac{\Big(\frac{d}{a} - \Big( \frac{b}{3a} \Big)^3 - \frac{b}{3a} \Big( \frac{c}{a} - 3 \Big( \frac{b}{3a} \Big)^2 \Big)^2}{4} +\frac{\Big( \frac{c}{a} - 3 \Big( \frac{b}{3a} \Big)^2 \Big)^3}{27}}}
  \\
  \intertext{Simplifying}
  c' = \frac{c}{a} - 3 \Big( \frac{b}{3a} \Big)^2 =\frac{3ac - b^2}{3 a^2}
  \\
  d' = \frac{d}{a} - \Big( \frac{b}{3a} \Big)^3 - \frac{b}{3a} c' = \frac{27 d a^2 + 2 b^3 - 9abc}{27 a^3}
  \\
  z = \sqrt[3]{-\frac{27 d a^2 + 2 b^3 - 9abc}{54 a^3} + \sqrt[2]{\frac{\Big( 27 d a^2 + 2 b^3 - 9abc \Big)^2}{54^2 a^6} +  \frac{\Big (3ac - b^2 \Big)^3}{3^6 a^6}   }}
  \\
  +
  \sqrt[3]{-\frac{27 d a^2 + 2 b^3 - 9abc}{54 a^3} - \sqrt[2]{\frac{\Big( 27 d a^2 + 2 b^3 - 9abc \Big)^2}{54^2 a^6} + \frac{\Big (3ac - b^2 \Big)^3}{3^6 a^6}}}
  \\
  \intertext{Factoring out 3a}
  z = \sqrt[3]{-\frac{27 d a^2 + 2 b^3 - 9abc}{54 a^3} + \frac{1}{(3a)^3} \sqrt[2]{\frac{\Big( 27 d a^2 + 2 b^3 - 9abc \Big)^2}{2^2} + \Big (3ac - b^2 \Big)^3}}
  \\
  +
  \sqrt[3]{-\frac{27 d a^2 + 2 b^3 - 9abc}{54 a^3} -
    \frac{1}{(3a)^3} \sqrt[2]{\frac{\Big( 27 d a^2 + 2 b^3 - 9abc \Big)^2}{2^2} + \Big (3ac - b^2 \Big)^3}}
  \intertext{Simplifying $\frac{\Big( 27 d a^2 + 2 b^3 - 9abc \Big)^2}{2^2} + \Big (3ac - b^2 \Big)^3 = A + B$}
  A = \frac{27^2}{4} a^4 d^2
  + 27 a^2 b^3 d - \frac{27 \cdot 9}{2} a^3 bc d + b^6 - 9 a b^4 c + \frac{9^2}{4} a^2 b^2 c^2
  \\
  B = 3^3 a^3 c^3 - 27 a^2 b^2 c^2 + 9 a b^4 c - b^6
  \\
  A + B = \frac{27^2}{4} a^4 d^2
  + 27 a^2 b^3 d - \frac{27 \cdot 9}{2} a^3 bc d + 3^3 a^3 c^3 + \Big( \frac{9}{4} - 3 \Big)9 a^2 b^2 c^2
  \\
  A + B =
  \frac{27 a^2}{4}
  \Big(
  27 a^2 d^2 + 4 b^3 d - 18 abcd + 4 a c^3 - b^2 c^2
  \Big)
  \intertext{Note that this quantity is the negative of the discriminant of the cubic, denoted by $\Delta$.}
  \intertext{Plugging this back in we obtain}
  z = \sqrt[3]{-\frac{27 d a^2 + 2 b^3 - 9abc}{54 a^3} + \frac{1}{2 \cdot 3^2 a^2} \sqrt[2]{
      3 \Big(
      27 a^2 d^2 + 4 b^3 d - 18 abcd + 4 a c^3 - b^2 c^2
      \Big)
    }}
  \\
  +
  \sqrt[3]{-\frac{27 d a^2 + 2 b^3 - 9abc}{54 a^3} -
    \frac{1}{2 \cdot 3^2 a^2} \sqrt[2]{
      3
      \Big(
      27 a^2 d^2 + 4 b^3 d - 18 abcd + 4 a c^3 - b^2 c^2
      \Big)
    }}
  \\
  \intertext{Simplfying more}
  z =
  \frac{1}{\sqrt[3]{2} \cdot 3 a} \Bigg(
  \sqrt[3]{9abc - 27 d a^2 - 2 b^3  + 3a \sqrt[2]{
      3 \Big(
      27 a^2 d^2 + 4 b^3 d - 18 abcd + 4 a c^3 - b^2 c^2
      \Big)
    }}
  \\
  +
  \sqrt[3]{9abc - 27 d a^2 - 2 b^3 -
    3a \sqrt[2]{
      3
      \Big(
      27 a^2 d^2 + 4 b^3 d - 18 abcd + 4 a c^3 - b^2 c^2
      \Big)
    }}
  \Bigg)
  \\
  \intertext{Undoing the first shifting of $z=x + \frac{b}{3a}$}
  x = - \frac{b}{3a} +
  \frac{1}{\sqrt[3]{2} \cdot 3 a} \Bigg(
  \sqrt[3]{9abc - 27 d a^2 - 2 b^3  + 3a \sqrt[2]{
      3 \Big(
      27 a^2 d^2 + 4 b^3 d - 18 abcd + 4 a c^3 - b^2 c^2
      \Big)
    }}
  \\
  +
  \sqrt[3]{9abc - 27 d a^2 - 2 b^3 -
    3a \sqrt[2]{
      3
      \Big(
      27 a^2 d^2 + 4 b^3 d - 18 a b c d + 4 a c^3 - b^2 c^2
      \Big)
    }}
  \Bigg)
  \\
  \intertext{Writing the full solution in terms of the discriminant we get}
  x = - \frac{b}{3a} +
  \frac{1}{\sqrt[3]{2} \cdot 3 a} \Bigg(
  \sqrt[3]{9abc - 27 d a^2 - 2 b^3  + 3^{\frac{3}{2}}a \sqrt[2]{\Delta
    } i}
  +
  \sqrt[3]{9abc - 27 d a^2 - 2 b^3 -
    3^{\frac{3}{2}} a \sqrt[2]{\Delta} i }
  \Bigg)
\end{gather*}

\subsection{Discriminant}\label{eps:subsec:discriminant2}
The discriminant of a cubic equation $a x^3 + b x^2 + c x + d = 0$ is defined to be
\begin{equation}
  \Delta = 18a b c d - 4b^{3}d + b^2 c^2 - 4ac^3 - 27a^2 d^2 \label{eps:eq:cubic-discriminant}
\end{equation}

\begin{enumerate}
  \item If $\Delta > 0$, the cubic has 3 distinct real roots.
  \item If $\Delta = 0$, the cubic has multiple roots (and all roots are real).
  \item If $\Delta < 0$, the cubic has 1 real root and 2 complex conjugate roots.
\end{enumerate}


\section{Quartic Equations}\label{eps:sec:quartic}

A general quartic polynomial has the form:
\begin{equation}
    a x^4 + b x^3 + c x^2 + d x + e = 0, \qquad (a \ne 0)
\end{equation}

\subsection{Depressing the Quartic}
Dividing by $a$ and applying the Tschirnhaus transformation $x = y - \frac{b}{4a}$ eliminates the cubic term, yielding the \emph{depressed quartic}:
\begin{equation}
    y^4 + p y^2 + q y + r = 0 \label{eps:eq:depressed-quartic}
\end{equation}
where
\begin{align*}
    p &= \frac{8ac - 3b^2}{8a^2}, \\
    q &= \frac{b^3 - 4abc + 8a^2d}{8a^3}, \\
    r &= \frac{-3b^4 + 256a^3e - 64a^2bd + 16ab^2c}{256a^4}
\end{align*}

\subsection{Ferrari's Method of Solution}
Rewrite equation~\eqref{eps:eq:depressed-quartic} by moving the linear and constant terms to the right-hand side:
\[
    y^4 + p y^2 = -q y - r
\]
Complete the square on the left-hand side by adding $\left(\frac{p}{2}\right)^2$:
\[
    \left(y^2 + \frac{p}{2}\right)^2 = -q y - r + \frac{p^2}{4}
\]
Now introduce an auxiliary parameter $m$. Adding $2m\left(y^2 + \frac{p}{2}\right) + m^2$ to both sides yields:
\begin{equation}
    \left(y^2 + \frac{p}{2} + m\right)^2 = 2m y^2 - q y + \left(m^2 + pm + \frac{p^2}{4} - r\right) \label{eps:eq:ferrari-identity}
\end{equation}
The right-hand side is a quadratic in $y$: $A y^2 + B y + C$ where $A = 2m$, $B = -q$, and $C = m^2 + pm + \frac{p^2}{4} - r$.

For the right-hand side to be a perfect square $(k y + l)^2$, its discriminant with respect to $y$ must vanish ($B^2 - 4AC = 0$):
\[
    (-q)^2 - 4(2m)\left(m^2 + pm + \frac{p^2}{4} - r\right) = 0
\]
Expanding this condition gives \textbf{Ferrari's Resolvent Cubic}:
\begin{equation}
    \boxed{8m^3 + 8pm^2 + (2p^2 - 8r)m - q^2 = 0} \label{eps:eq:resolvent-cubic}
\end{equation}
This is a cubic equation in $m$, which can be solved explicitly in radicals using Cardano's formula.

Let $m_0 \ne 0$ be any real (or complex) root of the resolvent cubic. Then the right-hand side of equation~\eqref{eps:eq:ferrari-identity} factors as:
\[
    2m_0 y^2 - q y + \left(m_0^2 + pm_0 + \frac{p^2}{4} - r\right) = \left( \sqrt{2m_0} y - \frac{q}{2\sqrt{2m_0}} \right)^2
\]
Taking square roots on both sides gives two quadratic equations:
\begin{equation}
    y^2 + \frac{p}{2} + m_0 = \pm \left( \sqrt{2m_0} y - \frac{q}{2\sqrt{2m_0}} \right)
\end{equation}
Rearranging into standard quadratic form:
\begin{align}
    y^2 - \sqrt{2m_0} y + \left(\frac{p}{2} + m_0 + \frac{q}{2\sqrt{2m_0}}\right) &= 0 \\
    y^2 + \sqrt{2m_0} y + \left(\frac{p}{2} + m_0 - \frac{q}{2\sqrt{2m_0}}\right) &= 0
\end{align}
Solving each quadratic with the standard quadratic formula yields the four roots $y_1, y_2, y_3, y_4$, and substituting back $x_k = y_k - \frac{b}{4a}$ provides the exact solutions to the original quartic equation.

\begin{remark}[The Abel-Ruffini Theorem]
By the Abel-Ruffini Theorem (1799/1824) and Galois Theory, the quartic is the highest-degree general polynomial equation solvable in radicals; general polynomials of degree $n \ge 5$ cannot be solved in terms of radicals.
\end{remark}


\section{Sylvester Matrices, Resultants, and General Discriminants}\label{eps:sec:sylvester-resultants}

\subsection{The Sylvester Matrix}

Let $P(x)$ and $Q(x)$ be polynomials with degrees $n \ge 1$ and $m \ge 1$:
\begin{align}
    P(x) &= a_n x^n + a_{n-1} x^{n-1} + \dots + a_1 x + a_0, \qquad (a_n \ne 0) \\
    Q(x) &= b_m x^m + b_{m-1} x^{m-1} + \dots + b_1 x + b_0, \qquad (b_m \ne 0)
\end{align}

\begin{definition}[Sylvester Matrix]\label{eps:def:sylvester-matrix}
The \textbf{Sylvester matrix} $S(P, Q)$ is the $(n+m) \times (n+m)$ matrix formed by filling $m$ rows with the coefficients of $P(x)$ and $n$ rows with the coefficients of $Q(x)$, shifted successively:
\begin{equation}
    S(P, Q) = 
    \begin{pmatrix}
        a_n & a_{n-1} & \cdots & a_1 & a_0 & 0 & \cdots & 0 \\
        0 & a_n & a_{n-1} & \cdots & a_1 & a_0 & \cdots & 0 \\
        \vdots & & \ddots & & & & \ddots & \vdots \\
        0 & \cdots & 0 & a_n & a_{n-1} & \cdots & a_1 & a_0 \\
        b_m & b_{m-1} & \cdots & b_1 & b_0 & 0 & \cdots & 0 \\
        0 & b_m & b_{m-1} & \cdots & b_1 & b_0 & \cdots & 0 \\
        \vdots & & \ddots & & & & \ddots & \vdots \\
        0 & \cdots & 0 & b_m & b_{m-1} & \cdots & b_1 & b_0
    \end{pmatrix}
\end{equation}
\end{definition}

\subsection{The Resultant}

\begin{definition}[Resultant]\label{eps:def:resultant}
The \textbf{resultant} of polynomials $P$ and $Q$, denoted $\operatorname{Res}(P, Q)$, is the determinant of their Sylvester matrix:
\begin{equation}
    \operatorname{Res}(P, Q) = \det S(P, Q)
\end{equation}
\end{definition}

\begin{theorem}[Fundamental Property of Resultants]\label{eps:thm:resultant-roots}
Let $P(x) = a_n \prod_{i=1}^n (x - \alpha_i)$ and $Q(x) = b_m \prod_{j=1}^m (x - \beta_j)$. Then:
\begin{equation}
    \operatorname{Res}(P, Q) = a_n^m b_m^n \prod_{i=1}^n \prod_{j=1}^m (\alpha_i - \beta_j) = a_n^m \prod_{i=1}^n Q(\alpha_i) = (-1)^{nm} b_m^n \prod_{j=1}^m P(\beta_j)
\end{equation}
In particular, $\operatorname{Res}(P, Q) = 0$ if and only if $P(x)$ and $Q(x)$ share a common root in $\mathbb{C}$.
\end{theorem}

\subsection{The General Polynomial Discriminant}

A polynomial $P(x)$ has a multiple root if and only if $P(x)$ and its derivative $P'(x)$ share a root, which occurs precisely when $\operatorname{Res}(P, P') = 0$.

\begin{definition}[Polynomial Discriminant]\label{eps:def:polynomial-discriminant}
For a polynomial $P(x) = a_n \prod_{i=1}^n (x - \alpha_i)$, the \textbf{discriminant} $\Delta(P)$ is defined as:
\begin{equation}
    \Delta(P) = a_n^{2n-2} \prod_{1 \le i < j \le n} (\alpha_i - \alpha_j)^2
\end{equation}
In terms of the resultant with its derivative $P'(x)$:
\begin{equation}
    \boxed{\Delta(P) = \frac{(-1)^{\frac{n(n-1)}{2}}}{a_n} \operatorname{Res}(P, P')}
\end{equation}
\end{definition}

\subsubsection{Example 1: Quadratic Discriminant}
For $P(x) = a x^2 + b x + c$, $P'(x) = 2a x + b$. The Sylvester matrix is $3 \times 3$:
\[
    S(P, P') = 
    \begin{pmatrix}
        a & b & c \\
        2a & b & 0 \\
        0 & 2a & b
    \end{pmatrix}
\]
Evaluating the determinant:
\[
    \det S(P, P') = a(b^2) - b(2ab) + c(4a^2) = a b^2 - 2a b^2 + 4a^2 c = -a(b^2 - 4ac)
\]
Since $n=2$, $(-1)^{n(n-1)/2} = (-1)^1 = -1$, yielding:
\[
    \Delta(P) = \frac{-1}{a} \left[ -a(b^2 - 4ac) \right] = b^2 - 4ac
\]

\subsubsection{Example 2: Cubic Discriminant}
For $P(x) = a x^3 + b x^2 + c x + d$, $P'(x) = 3a x^2 + 2b x + c$. The Sylvester matrix is $5 \times 5$:
\[
    S(P, P') = 
    \begin{pmatrix}
        a & b & c & d & 0 \\
        0 & a & b & c & d \\
        3a & 2b & c & 0 & 0 \\
        0 & 3a & 2b & c & 0 \\
        0 & 0 & 3a & 2b & c
    \end{pmatrix}
\]
Since $n=3$, $(-1)^{n(n-1)/2} = (-1)^3 = -1$. Evaluating $\det S(P, P')$ directly gives:
\[
    \Delta(P) = -\frac{1}{a} \det S(P, P') = 18abcd - 4b^3d + b^2c^2 - 4ac^3 - 27a^2d^2
\]
matching exactly the formula in equation~\eqref{eps:eq:cubic-discriminant}.

\end{document}
